% Unified 13-Tuple Formalism

Signal hierarchy theory (Section 3) introduces signal places ($\Psi$), signal flow arcs with consumption semantics, hierarchical layers ($L: \Psi \to \mathbb{N}$), and preemption mechanisms. We now integrate these concepts into a unified formalism extending classical Bio-PN with formal semantics for hierarchical regulatory control. The resulting Signal Hierarchical Petri Net (SHYPN) provides complete mathematical specification for modeling integrated metabolic-regulatory systems.

\subsection{Complete Definition}

\begin{definition}[Signal Hierarchical Petri Net]
A Signal Hierarchical Petri Net (SHYPN) is a 13-tuple:
\begin{align*}
\text{PN} = \langle &P, T, \text{Pre}, \text{Post}, m_0, k, S, \\
&\Phi, \Sigma, \text{Reg}, \Psi, E, A \rangle
\end{align*}
\end{definition}

\textbf{Component descriptions}:

\begin{table}[h]
\caption{13-tuple SHYPN components}
\centering
\small
\begin{tabular}{@{}cp{1.0cm}p{5.8cm}@{}}
\toprule
\# & \textbf{Tuple} & \textbf{Description} \\
\midrule
1 & $P$ & Places (biochemical species) \\
2 & $T$ & Transitions (reactions) \\
3 & Pre & Input arcs: $P \times T \to \mathbb{R}_{\geq 0}$ \\
4 & Post & Output arcs: $T \times P \to \mathbb{R}_{\geq 0}$ \\
5 & $m_0$ & Initial marking: $P \to \mathbb{R}_{\geq 0}$ \\
6 & $k$ & Rate constants: $T \to \mathbb{R}_{> 0}$ \\
7 & $S$ & Transition types \\
8 & $\Phi$ & Rate functions: $T \to (\mathbb{R}^{|P|} \to \mathbb{R}_{\geq 0})$ \\
9 & $\Sigma$ & Regulatory places: $\Sigma \subseteq P$ \\
10 & Reg & Regulatory modulation: $T \times \Sigma \to \mathbb{R}$ \\
11 & $\Psi$ & Signal places: $\Psi \subseteq P$ \textit{(new)} \\
12 & $E$ & Signal thresholds: $T \times \Psi \to \mathbb{R}_{\geq 0}$ \textit{(new)} \\
13 & $A$ & Arc types \textit{(new)} \\
\bottomrule
\end{tabular}
\end{table}

\vspace{3pt}
\noindent\textbf{Extensions beyond classical Bio-PN}~\cite{heiner2008}: \textbf{(1) $\Psi$}---Formal signal places for hierarchical information channels \textbf{(2) $E$}---Enablement thresholds for signal flow arcs \textbf{(3) $A$}---Arc type classification including novel signal\_flow type

\subsection{Arc Type Taxonomy}

\vspace{3pt}
\noindent Four arc types with distinct semantics:

\vspace{3pt}
\noindent\textbf{(1) Normal} ($A(a) = \text{normal}$): Mass transfer---Consumes source tokens: $M(p) \gets M(p) - \text{Pre}(p, t)$---Produces target tokens: $M(p) \gets M(p) + \text{Post}(t, p)$---Stoichiometric weights define conservation

\vspace{3pt}
\noindent\textbf{(2) Test} ($A(a) = \text{test}$): Catalytic read---Does \textbf{not} consume: $M(p)$ unchanged---Enables transition: $M(p) \geq \text{threshold}$---Models enzymes, catalysts, cofactors

\vspace{3pt}
\noindent\textbf{(3) Inhibitor} ($A(a) = \text{inhibitor}$): Allosteric inhibition---Does \textbf{not} consume: $M(p)$ unchanged---Disables transition: $M(p) < \text{threshold}$---Models competitive/non-competitive inhibition

\vspace{3pt}
\noindent\textbf{(4) Signal\_flow} ($A(a) = \text{signal\_flow}$): Information transfer \textit{(new)}---\textbf{Does} consume: $M(p_{\text{signal}}) \gets M(p_{\text{signal}}) - W(a)$---Hierarchical control: Layer-dependent execution---Models regulatory signals, checkpoints, decision finality

\subsection{Dynamics}

\vspace{3pt}
\noindent\textbf{Marking evolution}: For continuous semantics:
\begin{align*}
\frac{dM(p)}{dt} = &\sum_{t \in {}^\bullet p} \Phi(t) \cdot \text{Post}(t, p) \\
&- \sum_{t \in p^\bullet} \Phi(t) \cdot \text{Pre}(p, t)
\end{align*}

\vspace{3pt}
\noindent\textbf{Rate function}: Integrates dependencies:
\begin{align*}
\Phi(t) = k_t &\cdot f_{\text{sub}}({}^\bullet t) \\
&\cdot f_{\text{reg}}(\Sigma(t)) \\
&\cdot f_{\text{sig}}(\Psi(t))
\end{align*}

\vspace{3pt}
\noindent\textbf{Enabling condition}: Transition $t$ enabled if: \textbf{(1) Substrates available}---$\forall p \in {}^\bullet t: M(p) \geq \text{Pre}(p, t)$ \textbf{(2) Regulatory conditions met}---Test arc thresholds satisfied, inhibitor arcs below threshold \textbf{(3) Signal conditions met}---$\forall p \in \Psi(t): M(p) \geq E(t, p)$

\subsection{Integration with Signal Hierarchy}

\vspace{3pt}
\noindent The unified formalism provides hierarchical execution:

\vspace{3pt}
\noindent\textbf{Signal hierarchy (Vertical dimension)}---Signal places $\Psi$ organized in layers via assignment $L: \Psi \to \mathbb{N}$---Higher layers execute first (priority-based scheduling)---Signal consumption preempts lower-layer alternatives, collapsing decision spaces

\vspace{3pt}
\noindent\textbf{Metabolic concurrency (Horizontal dimension)}---Multiple metabolic reactions with disjoint substrates execute concurrently---Rate contributions accumulate for shared metabolic products (e.g., ATP from glycolysis and oxidative phosphorylation)

\vspace{3pt}
\noindent\textbf{Combined architecture}:
\begin{center}
\begin{verbatim}
L2: [Sig Psi_2] -> hier. ctrl
       | (signal flow)
L1: [Sig Psi_1] -> effector
       |
L0: [T_1||T_2||T_3] -> reactions
     |    |    |
    [P_1][P_2][P_3] -> pools
\end{verbatim}
\end{center}

\vspace{3pt}
\noindent This 2D architecture is unique in Petri net literature. Recent work addresses resource constraints~\cite{genovese2021} but not hierarchical integration.
